% Prepared for Edward B. Baker III with OpenAI Codex assistance, 26 November 2026.
% Consolidates the checked-in surgery notes, peripheral and round control,
% and independent reviews at ece209c. Scientific input freeze c8fb08dc.
% Research lead selection GPT-6.2 Astra/ultra. Bounded derivations and readers
% were dispatched as GPT-6.2 Sol/ultra. Actual serving model and effort were
% unavailable for independent verification, including for this editorial pass.
\documentclass[11pt,letterpaper]{article}
\usepackage[margin=1in]{geometry}
\usepackage[T1]{fontenc}
\usepackage{lmodern,microtype,amsmath,amssymb,amsthm,mathtools,booktabs}
\usepackage[hidelinks]{hyperref}
\setlength{\emergencystretch}{2em}
\numberwithin{equation}{section}

\newtheorem{theorem}{Theorem}[section]
\newtheorem{proposition}[theorem]{Proposition}
\newtheorem{lemma}[theorem]{Lemma}
\newtheorem{corollary}[theorem]{Corollary}
\theoremstyle{remark}
\newtheorem{example}[theorem]{Example}
\newcommand{\HB}[1]{\mathcal H_{#1}}
\newcommand{\TB}[1]{\mathcal T_{#1}}
\newcommand{\rd}{\gamma_{\mathrm{rd}}}
\DeclareMathOperator{\tr}{tr}
\DeclareMathOperator{\divg}{div}
\DeclareMathOperator{\Ric}{Ric}
\DeclareMathOperator{\intt}{int}

\title{Smooth recognition from marked handle cuts}
\author{Edward B. Baker III}
\date{26 November 2026}

\begin{document}
\maketitle

\begin{abstract}
Consider a smooth homotopy four-sphere with one lower and one upper
odd-index handle and two middle handles. After framed slides normalize
the lower winding row to \((1,0)\). A nonzero-degree map from a closed
oriented three-manifold of positive scalar curvature identifies the
winding-one intermediate cut as \(S^3\), or the winding-zero cut as
\(\#_2(S^1\times S^2)\), according to the order tested. Either
identification recognizes the original smooth sphere. We give
quantitative scalar curvature criteria for the actual marked cuts and
a cancellation of a certified winding-one attachment relative to any
finite number of movable framed tubes. For three middle handles a
second certified prefix recognizes the resulting two-component
surgery trace. An independent finite comparison on a subcritical
minimizing sequence shows that an exactly conformally round, umbilic
face of a parametrized disk forces equal chiral disk masses bounded by
the minimizing excess. We derive the required full-side savings for
balanced enlargements with a protected core.
\end{abstract}

\section{The marked surgery model}\label{sec:mark}
Let \(M\) be a smooth, closed, connected, oriented homotopy
four-sphere. Any metric used below is smooth and Riemannian. Give
each regular level the outward boundary orientation from the lower
sublevel. Reduce the lower and upper dual 0/1 handle graphs along
spanning trees and order the handles by index. If both ends have a
single extremum, the resulting counts are
\[
 (1,\ell,n_2,r,1),\qquad \ell+r=n_2.
\]
The last identity follows from \(\chi(M)=2\). The two end bodies
are \(P\cong\HB{\ell}\) and \(Q\cong\HB{r}\), where
\[
 \HB{j}=\natural_j(S^1\times D^3),\qquad
 \TB{j}=\natural_j(S^2\times D^2),\qquad
 Y_j=\#_j(S^1\times S^2).
\]
Here \(\natural\) is boundary connected sum and \(Y_0=S^3\).
In the mixed case \(\ell=r=1\). The tree reductions transfer
the actual handles and end identifications
\cite[Lemma 4.7 and Theorem 4.8]{Milnor1965HCobordism}
\cite[marked end reductions]{Baker2026MarkedHeights}.

Write \(N\) for the ordered middle cobordism, \(p:\partial P\to
\partial_-N\) for its lower gluing, and \(q:\partial Q\to\partial_+N\)
for its upper gluing. If \(\widetilde a_j:S^1\times D^2\to\partial_-N\)
is an incoming tube, use \(a_j=p^{-1}\widetilde a_j\) on \(\partial P\).
For a subset \(I\) of the disjoint 2-handle attachments let \(X_I\) be
the lower body with those attachments and write \(Z_I=\partial X_I\).
The surgery coordinate model replaces each tube for \(j\in I\)
by its outgoing face \(b_j:D^2\times S^1\to Z_I\). On the common torus
the map and inverse identify
\[
 b_j(z,w)\longleftrightarrow a_j(z,w),\qquad |z|=|w|=1.
\]
The exterior is carried by the chosen product collars, with reverse
flow for the inverse. An unfilled tube retains its full map \(a_j\)
read in that exterior. Compatible collars and roundings therefore
identify the outputs of all orders at the full subset \(F\). This
identification is the identity on the common exterior and the
parametrized outgoing solid tori. Both product coordinates of a
2-handle are required in this comparison
\cite[quadratic handle model]{Baker2026ImmersedHandles}.

For the resulting outgoing mark \(\eta:Z_F\to\partial_+N\), the
closed assembly uses \(\varphi=q^{-1}\eta:Z_F\to\partial Q\). It
reverses the two outward orientations. Read downwards, the
attaching tube of dual handle \(j\) is
\begin{equation}\label{eq:dual}
 a_j^*(w,z)=\varphi b_j(z,w),\qquad
 (a_j^*)^{-1}(y)=\operatorname{swap}\bigl(b_j^{-1}\varphi^{-1}(y)\bigr).
\end{equation}
The \(\operatorname{swap}\) exchanges the two coordinates of the
product handle. Equation~\eqref{eq:dual} uses the initial top
coordinates for the first reverse attachment. Subsequent dual
tubes use the partial-level marking transported through each
preceding surgery on the surviving exterior, with its inverse.
Each reverse replacement recovers the original tube \(a_j\).
After all reverse surgeries the lower comparison is \(p^{-1}\).
Changing the order of the forward attachments and reversing the
height are distinct operations.

A comparison \(F_X:X_F'\to X_F\) of lower traces, including a framed
handle slide, returns the upper maps
\begin{equation}\label{eq:reglue}
 \varphi'=\varphi\,\partial F_X,\qquad
 (\varphi')^{-1}=(\partial F_X)^{-1}\varphi^{-1}.
\end{equation}
The closed comparison uses \(F_X\) and the identity on \(Q\).
When \(F_X\) restricts to \(F_P\) on the model lower body, its lower
gluing is \(p\partial F_P\) with inverse
\((\partial F_P)^{-1}p^{-1}\). These rules carry framings and
exterior collars as part of the smooth identification. A metric
on a particular level of a particular source requires the
embedded comparison with that level.

\section{A control with two framed attachments}\label{sec:control}
Use \(P=S^1_z\times D^2_w\times[0,1]\) with rounded corners and
lower belt sphere \(B=\{1\}\times\partial(D^2_w\times[0,1])\).
Attach \(K\) by the product tube \(a_K(z,w)=(z,w,1)\).
The core crosses \(B\) once with positive sign and represents
the zero-framed meridian of a dotted 1-handle. Place a second
component \(S\) with its disk framing in a small ball in the
untouched boundary disjoint from \(a_K\) and \(B\). At \(K\)
first the lower ball map and its inverse are
\begin{equation}\label{eq:radial}
 \begin{aligned}
 G_K:P\cup h_K&\longrightarrow D^2_{|Z|\le2}\times D^2_w,\\
 (z,w,t)&\longmapsto ((2-t)z,w),&
 (Z,w)_{h_K}&\longmapsto (Z,w).
 \end{aligned}
\end{equation}
The inverse uses the handle for \(|Z|\le1\), and
\((Z/|Z|,w,2-|Z|)\) for \(|Z|\ge1\). The signed collar
coordinate at the gluing is \(|Z|-1\). Round the remaining
corners consistently in these coordinates. The transported
spectator tube is \((\partial G_K)a_S\) in the same exterior.

To mark the spectator surgery, identify the boundary of the
rounded lower ball with the unit \(S^3\subset\mathbb C^2\)
by an angular map \(L\) with
\[
 L a_S(z,v)=(\sqrt{1-\epsilon^2|v|^2}\,z,\epsilon v),
 \qquad |v|\le1,\quad 0<\epsilon<1.
\]
One can choose \(L\) by radial projection, followed by a rotation
and a stereographic dilation. Take the projection point off
the displayed tube and send a point of the untouched patch to
its antipode. The inverse dilation brings the compact displayed
tube into the patch for a sufficiently large dilation factor.
This constructs \(a_S\) by pullback and \(L^{-1}\) by reversing
the factors, including the reverse radial map. The exterior
of the displayed tube has coordinates
\begin{equation}\label{eq:two-disks}
 q_S=\frac{y}{|y|},\qquad u=\frac{x}{\sqrt{1-\epsilon^2}},
 \qquad
 (x,y)=\bigl(\sqrt{1-\epsilon^2}\,u,
       \sqrt{1-(1-\epsilon^2)|u|^2}\,q_S\bigr),
\end{equation}
for \((x,y)\in S^3\) with \(|y|\ge\epsilon\).
The new face \(b_S(u,q_S)\) agrees in both torus coordinates.
Thus the output is the double of the \(u\)-disk times
\(S^1_{q_S}\). Let \(\iota_-\) and \(\iota_+\) mark the exterior
and inserted disk of this double in \(S^2\). Choose the actual
closing map to \(Q=S^1_{q_S}\times D^3\) to be
\[
 \varphi(y_0)=(q_S,\iota_-(u)),\qquad
 \varphi b_S(u,q_S)=(q_S,\iota_+(u)).
\]
The first formula uses \eqref{eq:two-disks} on \(L(y_0)\). Its
inverse uses the displayed \((x,y)\) followed by \(L^{-1}\),
or \(b_S(u,q_S)\) on the inserted hemisphere. Give \(Q\) the
orientation for which \(\varphi\) reverses outward boundary
orientations. The upper product sphere \(\{q_S=1\}\times S^2\) meets the spectator belt once.
With \(\iota_+\) the top disk of the upper \(D^3\) presented
as a product, \eqref{eq:radial} in the \(q_S\) coordinate
gives the upper ball map from \(Q\cup h_S^*\) and its inverse.
The two complete sides at \(Z_{\{K\}}\) are balls. Their
actual gluing gives the standard \(S^4\), since every boundary
diffeomorphism of \(S^3\) extends over \(D^4\)
\cite[p.~553]{Hatcher1983Smale}.

If \(S\) is attached first, its local zero-framed surgery
replaces its ball by a punctured \(S^1\times S^2\) via the
two-disk marking. The resulting cut is \(Y_2\). The comparison
is the punctured version of \eqref{eq:two-disks} inside the
ball and the identity outside, likewise for its inverse.
The remaining \(K\) tube is the original tube in that exterior.

\subsection{A framed slide and the dual cuts}
Use a slightly enlarged \(K\) chart to take a product push-off
\(K^+\) just outside \(a_K\). Choose a narrow embedded band from
\(S\) to \(K^+\) in the complement of the belt. Its interior
misses the thick tubes away from its prescribed endpoint
collars, and its split-disk endpoint approaches the disk edge
from the exterior. Use the untwisted thickening with the
prescribed parallels at both ends. This gives the framed
band sum \(S'\) and changes the oriented winding row from
\((1,0)\) to \((1,1)\)
\cite[\S2 and Lemma 2.1]{GompfScharlemannThompson2010}.
In the \(K\) surgery the push-off bounds the disk
\(D^2_z\times\{w_0\}\) in its outgoing face, extended by a
short exterior annulus. Nearby values of \(w_0\) give
parallel disks realizing the framing. Shrinking the band
detour across this disk gives a framed isotopy after \(K\).
Write \(J_t\) for an ambient extension, parametrized so
\(J_1a_S=a_{S'}\) in the surviving exterior. Such an extension
is obtained by extending the velocity of the moving framed
tube with a cutoff and flowing the resulting vector field.

A new-to-old comparison \(F_X\) applies \(J_{\lambda(t)}^{-1}\)
in a pre-spectator collar of the output of \(K\), with
\(\lambda=0\) near the lower face and \(\lambda=1\) near
its upper face. It is the identity below this collar and in
the second handle coordinates. Its inverse uses
\(J_{\lambda(t)}\) on the same collar. Equation~\eqref{eq:reglue}
therefore identifies the closed smooth assemblies. The upper
\(S'\) dual tube agrees in the new marking with the previous
\(S\) dual tube since \(\partial F_X b_{S'}=b_S\). The other
dual tube is carried by \eqref{eq:dual} with \(\varphi'\). In handle
bases dual to the oriented forward basis
\((e_K,e_S+e_K)\) the dual basis is \((f_K-f_S,f_S)\).
The upper winding row in the slid coordinates is consequently
\((-1,1)\) if its previous row was \((0,1)\). Its signed
change agrees with the dual framed slide description.

For \(S'\) first, shrink the other band detour across the
split spanning disk of \(S\) in \(\partial P\) and translate
the parallel \(K^+\) to \(K\) in the product tube. The
ambient endpoint \(A\) can be taken with \(A a_{S'}=a_K\)
through the displayed untwisted disk framings. It transports
the residual \(a_K\) at the same time. Use \(A\) on the
surviving exterior and the identity on the replacing torus,
then apply \(\partial G_K\). This identifies the \(S'\)-first
cut with \(S^3\) and gives the framed residual tube there.
The inverse uses \((\partial G_K)^{-1}\) and \(A^{-1}\) on
the exterior, with the reverse torus identification. Extend the marked comparison
over the residual handle by the identity in its coordinates.
If \(\mathcal A\) is the resulting comparison from the slid trace
to the returned trace with the standard ball, its closing map is
\(\varphi'(\partial\mathcal A)^{-1}\).
The cut is \(S^3\) also when reached by attaching \(K^*\) first
from above. This spherical output certifies upper cancellation
after an extendible re-marking by Proposition~\ref{prop:peripheral},
with \(K^*\) oriented as winding one. The signed winding row
alone specifies only the algebraic count in the initial upper chart.
The intermediate cuts, including the first dual attachment
that reaches the same cut, are
\[
 \begin{array}{@{}llll@{}}
 \toprule
 \text{Diagram}&\text{First from below}&\text{First from above}
     &\text{Intermediate boundary}\\
 \midrule
 \text{Split}&K&S^*&S^3\\
 \text{Split}&S&K^*&Y_2\\
 \text{Slid}&K&(S')^*&S^3\\
 \text{Slid}&S'&K^*&S^3\\
 \bottomrule
 \end{array}
\]
The split and slid controls have the same middle count of two.
The indicated geometric lower or upper cancellation gives one
middle handle on the original type. At the \(K\)-first cut
both cancellations are available and the count becomes zero.
Milnor's cancellation condition is one
\emph{transverse geometric} intersection in the common level.
Under that condition his product flow and its inverse transport
arbitrary remaining attachments through the removal
\cite[Theorem 5.4]{Milnor1965HCobordism}.

\subsection{The full cluster and a metric on the spherical cut}
The full lower cluster is the lower 0/1 body with both entire
2-handles, all joining tubes and an outgoing collar, rounded at
the corners. In any mixed presentation
the actual upper gluing identifies its boundary with
\(\partial\HB{1}\cong S^1\times S^2\). This holds in either
order and after transported slides. Its nontrivial fundamental
group excludes this whole cluster as a spherical cut. A partial
neighborhood of the interaction or an intermediate section
has a different boundary identification.

At the split \(K\)-first cut use the two explicit ball maps to
choose a diffeomorphism \(D:M\to S^4\) extending \(L\) to the
two hemispheres. Hatcher's extension theorem adjusts each
boundary discrepancy over its ball. The maps can be collared
with opposite outward normals at the seam
\cite[p.~553]{Hatcher1983Smale}. Pull back the unit round metric
by \(D\). In the bicollar given by the inverse latitude
coordinates it is \(dq^2+\cos^2(q)\rd\). For two available
subwidths inside \(|q|<\pi/2\), choose a cutoff \(\beta\)
which is one near the equator and zero near both outer faces.
The conformal path
\[
 \exp(2s\beta(q)\log\sec q)\,D^*g_{S^4},
 \qquad 0\le s\le1 ,
\]
ends with an exact round product near the equator in the
coordinate \(t=\int_0^q\sec v\,dv\). Both complete chiral side
energies and the entire joined peak are zero. Every class
is positive, and the marked conformally flat caps have zero
energy and positive classes by the argument in
Section~\ref{sec:budget}. The same construction on the slid
control uses \(F_X\cup\mathrm{id}_Q\). It gives the physical
embedding, both complete sides and the inverse reassembly on
each control source.

\section{Recognition at intermediate levels}\label{sec:psc}
For a single lower 1-handle the oriented intersection numbers of
the incoming 2-handle cores with its belt sphere are their winding
numbers in \(S^1\times S^2\). The row is primitive since its cokernel
is \(H_1(M;\mathbb Z)=0\). Actual framed slides, changes of core
orientation and permutations implement the Euclidean algorithm.
Each slide carries the upper gluing by \eqref{eq:reglue}.
In the mixed case write the normalized row as \((1,0)\).
Let \(Z_a\) be the cut obtained with the winding-one attachment first,
and let \(Z_b\) be the cut in the other order.

Say a closed oriented three-manifold is \emph{PSC-dominated} when
a closed oriented three-manifold carrying a positive-scalar-curvature
metric maps to it with nonzero degree. A disconnected source suffices
if some component has nonzero degree. In particular a positive metric
on the target suffices by the identity map. Chodosh--Maximo--Mukherjee,
Proposition 2.2, gives the connected sum of spherical space forms
and \(S^1\times S^2\) factors for an orientable PSC-dominated target.
Corollary 2.3 specializes the decomposition on an integral homology
sphere to \(S^3\) or a connected sum of signed Poincar\'e homology spheres.
Their Proposition 4.2 rules out every nonempty such sum as the boundary
of a smooth compact contractible four-manifold
\cite[Proposition 2.2, Corollary 2.3 and Proposition 4.2]{ChodoshMaximoMukherjee2024Complete}.

\begin{lemma}[Zero-matrix surgery]\label{lem:zero}
Let \(V\) be the ordinary trace of integer-framed surgery on a link
of \(j\) components in \(S^3\) with zero linking matrix, including
the framings. If \(\partial V\) is PSC-dominated then
\(\partial V\cong Y_j\).
\end{lemma}
\begin{proof}
The trace is a ball with \(j\) 2-handles. Thus it is simply connected
and its intersection form on \(H_2(V)=\mathbb Z^j\) is zero.
The boundary first homology is the cokernel of the linking matrix.
The pair sequence and Poincar\'e--Lefschetz duality give
\[
 0=H^1(V)\cong H_3(V,\partial V)\longrightarrow H_2(\partial V)
 \longrightarrow H_2(V)\longrightarrow H_2(V,\partial V).
\]
The last arrow is the zero intersection map, so the inclusion in
the middle is an integral isomorphism. The scalar curvature
decomposition has \(j\) product factors. Each spherical summand
has trivial first homology since their finite first homology
groups form a direct sum in the torsion-free boundary group.
The remaining sum \(\Sigma\) of spherical summands has integral
homology of \(S^3\). Its spherical metrics and positive-scalar-curvature
connected sum supply a positive metric on \(\Sigma\)
\cite[Theorem A]{GromovLawson1980Surgery}.
Corollary 2.3 then identifies the nontrivial summands as signed
Poincar\'e spheres.

In this decomposition choose a transverse factor sphere in each
\(S^1\times S^2\) and pull the disjoint spheres back to the actual
boundary of \(V\). They give an integral basis of \(H_2(V)\) by
the pair sequence. Attach \(j\) auxiliary 3-handles along the
spheres with their product normal intervals. Compression of each
factor along its sphere and capping the two sides removes that
\(S^1\times S^2\) summand. The resulting smooth filling of
\(\Sigma\) is simply connected, since 3-handles preserve the
fundamental group. It is acyclic since the \(j\) relative
3-generators map isomorphically to \(H_2(V)\).
It is therefore contractible. Proposition 4.2 of
Chodosh--Maximo--Mukherjee gives \(\Sigma\cong S^3\).
\end{proof}

\begin{theorem}[The two normalized mixed cuts]\label{thm:mixed}
In a reduced mixed presentation of \(M\) with two middle handles,
\(Z_a\) is PSC-dominated if and only if \(Z_a\cong S^3\).
The cut \(Z_b\) is PSC-dominated if and only if \(Z_b\cong Y_2\).
Either condition implies that \(M\) is diffeomorphic to the standard
\(S^4\). The same statements apply to a normalized dual presentation
with its own marked cuts.
\end{theorem}
\begin{proof}
For \(Z_a\) the lower trace of the single attachment has a
winding-one relation killing the lower fundamental group
\(\mathbb Z\), and a unit cellular boundary in degree two.
It is smooth and contractible. Hence \(Z_a\) is an integral
homology sphere. The signed Poincar\'e sum obstruction just
stated gives \(Z_a\cong S^3\) under PSC domination.

For \(Z_b\) replace the lower dotted unknot \(U\) by a zero label
in the ordinary boundary surgery description. The full ordinary
surgery matrix is
\[
 \begin{pmatrix}0&1&0\\1&d&c\\0&c&e\end{pmatrix}
 \quad\text{in the order } U,K_a,K_b .
\]
Its determinant is \(-e\). The full boundary is marked by the
upper body \(\HB{1}\), so its first homology has rank one and
\(e=0\). Attaching \(K_b\) first gives the ordinary surgery
on \(U_0,K_b\), with zero matrix. Its trace from \(D^4\) is
a comparison filling of \(Z_b\). Lemma~\ref{lem:zero} applies
to that filling and gives \(Y_2\). The auxiliary conversion
of the dot and subsequent 3-handles identify the boundary
while the original lower trace retains its 1-handle.
Conversely the round metric on \(S^3\), and the connected sum
of two positive product metrics on \(S^1\times S^2\), give
positive scalar curvature. Here we use preservation of positive
scalar curvature under connected sum in dimension three
\cite[Theorem A]{GromovLawson1980Surgery}.

The marked mixed recognition theorem applies to either boundary
with the original attaching and closing maps
\cite[mixed recognition theorem]{Baker2026MarkedHeights}.
For \(S^3\), the cancellation and one-knot argument below also
prove recognition directly. To recall the geometric content of the
\(Y_2\) case, each half read from its own end is one dotted 1-handle
and one 2-handle with boundary \(Y_2\).
The proof of Gompf--Scharlemann--Thompson's Proposition 3.2
puts the latter component as a zero-framed unknot in a ball
in the actual lower \(S^1\times S^2\). It uses the ordinary
two-component link consisting of the dotted circle changed to
a zero-framed unknot and the given attachment. The isotopy in
the zero surgery on the unknot carries the framed tube
\cite[Proposition 3.2 and its proof]{GompfScharlemannThompson2010}.
Thus each half is \((S^1\times D^3)\natural(S^2\times D^2)\)
with its transported mark. The marked theorem uses interior
surgery on the \(S^2\) factor in one such half and the
1-handlebody boundary extension of Laudenbach--Po\'enaru to
give \(S^1\times S^3\). The dual circle of that surgery
represents a generator of its fundamental group since its
inverse surgery gives the simply connected \(M\). A homotopy
to a product generator gives an isotopy of circles by
parametric transversality in dimension four. Laudenbach--Po\'enaru's
Theorem A then handles the actual framing of the inverse
surgery. This is the \(Y_2\) argument of the marked theorem
\cite[Theorem A and its proof, pp.~337, 341--342]{LaudenbachPoenaru1972Handlebodies}
\cite[mixed recognition theorem]{Baker2026MarkedHeights}.
\end{proof}

\begin{example}[A published order test]
In the diagram of \(DW^-(0,n+2)\) in Isoshima's Figure 37,
the zero-framed small meridian of the long attaching component
has a ball disjoint from the dotted circle. Its spanning disk
may meet the long component that will be attached later.
The meridian-first surgery replaces that ball by punctured
\(S^1\times S^2\), so the cut is \(Y_2\) with the long framed
tube carried in the surviving exterior. Theorem~\ref{thm:mixed}
recognizes this marked closed diagram. At \(n=1\) the
long-component-first cut is, up to orientation,
\(\Sigma(2,5,7)\), with nontrivial fundamental group.
Its filling is the original Mazur piece. The same closed
diagram thus supplies a certified cut and a cut outside
the recognition condition for its order
\cite[Figure 37]{Isoshima2025MazurDoubles}
\cite[\S1, Theorem 2, and \S5, Figure 22]{AkbulutKirby1979Mazur}.
\end{example}

\section{Reducing a specified winding-one attachment}\label{sec:reduction}
The winding-one argument uses a single lower 1-handle and
one chosen attachment. The number of disjoint framed tubes
awaiting attachment can be any finite number.

\begin{proposition}[Relative cancellation and the peripheral test]\label{prop:peripheral}
Suppose a chosen winding-one component \(K\) is attached first
above a single lower 1-handle. Put
\(X=(S^1\times S^2)\setminus\intt\nu K\), and let \(\mu\)
be its tube meridian in \(\pi_1(X)\) up to conjugacy and inversion.
The following are equivalent:
\begin{enumerate}
 \item \(\mu\) is trivial in \(\pi_1(X)\).
 \item \(\pi_1(X)\cong\mathbb Z\).
 \item Surgery with the given framing of \(K\) yields \(S^3\).
 \item After an extendible re-marking of the lower body the core
       of the chosen tube meets the lower belt sphere once
       transversely.
\end{enumerate}
In any of these cases the lower pair cancels and \(n_2\) decreases
by exactly one on the original smooth type. The remaining framed
tubes and closing map are transported through the cancellation.
A PSC-dominated winding-one first cut satisfies the conditions.
\end{proposition}
\begin{proof}
If an exterior for the entire link is used, fill all spectator
tori in their original meridional slopes to recover this \(X\).
Filling the meridian of \(K\) quotients \(\pi_1(X)\) by the
normal closure of \(\mu\) and gives \(\mathbb Z\). This proves
the equivalence of the first two conditions. Framed surgery
also takes a quotient. The resulting group, if cyclic,
is trivial because the filled manifold has \(H_1=0\).
Indeed the dotted-to-zero two-by-two matrix of a winding-one
attachment has determinant \(-1\) for every integer framing.
The three-dimensional Poincar\'e theorem now gives \(S^3\)
\cite{Perelman2002,Perelman2003Surgery,Perelman2003Extinction}.

If the framed surgery yields \(S^3\), the dual knot there
has surgery \(S^1\times S^2\) by \eqref{eq:dual}. Its integer framing is
zero by the surgery homology computation. Gabai's Property R
gives the unknot. Its zero surgery has a product surgery core,
so \(K\) in the original boundary has that description in
the surgery coordinates. A boundary change between these
coordinates and the lower body extends over \(\HB{1}\).
Equivalently the reverse Property R theorem identifies this
core up to isotopy in the original boundary
\cite[Corollary 8.3 and Remark 8.5]{Gabai1987FoliationsIII}
\cite[Theorem 10.1]{Scharlemann1989Sutured}
\cite[proof of Theorem A, pp.~341--342]{LaudenbachPoenaru1972Handlebodies}.
The original normal parametrization is carried by this
comparison. Puncturing the re-marked belt sphere at its
single transverse intersection with the core gives a
meridian disk in the exterior after shrinking the tube.
This proves the first condition from the fourth.

Milnor's first cancellation theorem now gives a product
cobordism with its end collars. Use the forward and inverse
product flow to transport the other framed tubes and the
closing map. The count removes precisely the chosen 1/2 pair
\cite[Theorem 5.4]{Milnor1965HCobordism}.
The contractible filling argument for \(Z_a\) in
Theorem~\ref{thm:mixed} uses just the first attachment,
so it applies with any finite number of remaining tubes.
\end{proof}

When the original mixed presentation had two middle handles,
the cancellation leaves a single 2-handle on \(D^4\) whose
surgery is \(S^1\times S^2\) in the actual upper marking.
Again the framing is zero by homology and the knot unknotted
by Property R. Its trace is \(S^2\times D^2\) and
Laudenbach--Po\'enaru's Theorem A gives the standard \(S^4\)
for the transported closing map
\cite[Corollaries 8.3 and 8.6]{Gabai1987FoliationsIII}
\cite[Theorem A]{LaudenbachPoenaru1972Handlebodies}.

The peripheral element tests the chosen component after restoring
the spectator tori. A nonabelian quotient of this restored
exterior group obstructs the direct cancellation of that component.
The obstruction pertains to this fixed cut. Actual slides that
change the chosen knot or a different ordering can produce a
different test. One may also probe the winding-one component
with a different integer framing. Every probe trace has a
contractible filling, so a PSC-dominated probe gives \(S^3\)
and identifies the product core. A Heegaard Floer \(L\)-space
probe of such a Mazur trace also identifies the product core
by Hayden--Mark--Piccirillo, Proposition 3.2. In both probes
the resulting knot identification is applied to the original
normal parametrization and to the original closed assembly
\cite[Proposition 3.2]{HaydenMarkPiccirillo2021Mazur}.

\subsection{A relative bigon}
For three transverse crossings of the chosen core with signs
\(+,+,-\) up to reversal, suppose a subarc between consecutive
opposite crossings and an arc on the belt bound an embedded disk.
Require the interiors of the arcs to be disjoint and the disk
interior to miss the chosen core and the belt. Assume the one-sided
disk and belt patch have a product neighborhood after corner
rounding. Isotope the first arc across the disk to a parallel of
the belt arc on the other side. The opposite signs allow both end
segments to stay on that side, so precisely the two specified
crossings disappear. Support the movement in the product
neighborhood, away from the third crossing and from the rest of
the chosen core. Extend the moving tube to an ambient isotopy
\(H_t\) of the common level. Use the belt as a fixed reference
for the new diagram and apply \(H_1\) to every framed tube.
The ambient diffeomorphism preserves disjointness of the tubes
and allows the disk sweep to meet spectators. It preserves their
signed winding numbers and can change their geometric crossing
counts. A spectator fixed in place requires a support disjointness
condition.

The closed comparison uses \(H_{\lambda(t)}^{-1}\) in a
pre-attachment product collar and the identity in the handle
coordinates, with inverse \(H_{\lambda(t)}\).
Equation~\eqref{eq:reglue} therefore carries the actual upper
assembly. After the move the core crosses once, so
Proposition~\ref{prop:peripheral} applies.
For a local example take \(B=\{z=0\}\) and the chosen arc
\(y=0,\ z=x(x^2-a^2)\), \(a>0\) small. The bigon spans
\([-a,0]\) in the half space \(z\ge0\).
At \(x=-a/2\) a disk-framed circle in the \(yz\)-plane
centered at the origin with radius \(a^3/8\) meets the
bigon and misses the chosen arc. It crosses the belt twice.
A cutoff translation of \(z\) by \(-a^3\), supported in
\(x\in(-a-\delta,\delta)\), \(0<\delta<a/4\), and constant
on the swept graph and spectator disk removes the first two
crossings. The third remains at \(x=a\). The circle and its
disk move to the negative side. The inverse motion realizes
this local configuration inside a split control with its
whole framing. In a general three-crossing diagram a disk
with the stated properties still has to be selected.

\subsection{Two prefix tests for three middle handles}
\begin{theorem}\label{thm:three}
Take a reduced \((1,1,3,2,1)\) presentation and normalize
the row by marked slides to \((1,0,0)\). For the order
\(a,b,c\) with \(a\) of winding one, let \(Z_a,Z_{ab}\)
be the first two actual boundaries. If both are
PSC-dominated, the original smooth sphere is standard.
If the first is PSC-dominated, the original sphere has a
presentation with two middle handles and no lower
1-handle. Apply in whichever direction has a single lower
1-handle. An original \((1,2,3,1,1)\) presentation qualifies
after reversal with its own cuts and normalization.
\end{theorem}
\begin{proof}
The first boundary is \(S^3\) by the winding-one
contractible filling and the obstruction to the signed
Poincar\'e sum. Proposition~\ref{prop:peripheral} cancels
the lower pair with both residual tubes present.
Let \(\psi:D^4\to X_a\) be the ball comparison returned
by the product cancellation, and let \(u_b,u_c\) denote
the two actual residual tubes on \(Z_a\). Their new
embeddings are \(v_j=(\partial\psi)^{-1}u_j\).
Let \(V_b\) and \(V_{bc}\) be the corresponding traces
from the ball. The comparisons
\(\Psi_b:V_b\to X_{ab}\) and \(\Psi_{bc}:V_{bc}\to X_{abc}\)
use \(\psi\) and the identity in the attached handle
coordinates. In particular the new top map is
\(\varphi\partial\Psi_{bc}\).
A map certifying \(Z_{ab}\) is transported to the
boundary of \(V_b\) by \((\partial\Psi_b)^{-1}\) and
keeps nonzero degree.

The entire surgery on \(v_b,v_c\) has boundary \(Y_2\)
by the actual upper \(\HB{2}\). The linking matrix is
zero since its integral cokernel is \(\mathbb Z^2\).
Thus \(V_b\) is a zero-framed one-knot trace.
Lemma~\ref{lem:zero} with \(j=1\) gives \(Y_1\)
for the PSC-dominated boundary of \(V_b\). Property R
makes \(v_b\) an unknotted component of the actual
remaining two-component link
\cite[Corollary 8.3]{Gabai1987FoliationsIII}.
Gompf--Scharlemann--Thompson's Proposition 3.2
now slides the actual link to the zero-framed unlink.
Let \(\Theta:\TB{2}\to V_{bc}\) be the resulting trace
comparison. The final gluing is
\(\varphi\,\partial\Psi_{bc}\,\partial\Theta\).
Laudenbach--Po\'enaru's Theorem A for rank two
applies to this gluing
\cite[Proposition 3.2]{GompfScharlemannThompson2010}
\cite[Theorem A]{LaudenbachPoenaru1972Handlebodies}.
\end{proof}

The first prefix gives a strict \(3\to2\) decrease before
the second prefix is tested. The output is a one-sided
two-component presentation, to which the separate
tunnel-number and generalized-square-knot recognition tests
can also apply under their respective hypotheses
\cite[one-sided recognition and stabilization]{Baker2026MarkedHeights}
\cite{MeierSchirmerZupan2016Trisections,MeierZupan2022SquareKnots}.
For more residual attachments the same first cancellation
returns a smaller presentation. Consider reduced counts
\((1,1,n_2,n_2-1,1)\) and the normalized winding row
\((1,0,\ldots,0)\). If a winding-zero component is tested first,
its ordinary two-component surgery with the dotted unknot
has zero matrix. Indeed the full matrix has rank two by the
actual upper boundary \(Y_{n_2-1}\).
The first two rows and columns give the block
\(\left(\begin{smallmatrix}0&1\\1&d\end{smallmatrix}\right)\).
Its inverse has bottom-right entry zero, so the Schur
complement is the original winding-zero linking block,
which must vanish. Lemma~\ref{lem:zero} identifies a
PSC-dominated winding-zero cut with \(Y_2\).
The proof of Proposition 3.2 then puts that component into
a ball with disk framing, carrying the other tubes. A net
2/3 cancellation in the longer presentation additionally
requires the upper attaching data.

\section{Metric criteria on the given levels}\label{sec:bands}
These criteria give PSC domination for the actual slice
or prefix in Theorems~\ref{thm:mixed} and~\ref{thm:three}.
Define \(\rho=\min_C R_g>0\) on a smooth compact marked
product band \(C=Z\times[t_-,t_+]\). Its width \(d\) is
the intrinsic Riemannian distance inside \(C\) between its ends.

\begin{proposition}[Scalar and end-mean bounds]\label{prop:band}
If the target \(Z\) has no positive-scalar-curvature metric,
then
\begin{align}
 \rho d^2&<3\pi^2,\label{eq:width}\\
 \arctan\frac{2a_-}{\sqrt{3\rho}}+
 \arctan\frac{2a_+}{\sqrt{3\rho}}+
 \sqrt{\rho/3}\,d&\le0,\label{eq:mean}
\end{align}
where \(a_\pm\) are the infima of the two boundary means
\(\divg_{\partial C}\nu_{\rm out}\). Hence the inclusive
width condition \(\rho d^2\ge3\pi^2\), or strict positivity
of the left side of \eqref{eq:mean}, certifies the cut.
\end{proposition}
\begin{proof}
Every closed separating hypersurface in the product has an
oriented collection of components homologous to the lower
slice. The sum of their degrees on projection to \(Z\) is
one. A positive metric on such a separator would give a
PSC-dominating component and therefore a positive metric on
\(Z\). Thus the band has the no-positive-separator property
required for R\"ade's comparison. His Corollary 2.25 on
\((\rho/12)g\) in dimension four gives
\(\sqrt{\rho/12}\,d<\pi/2\), proving \eqref{eq:width}.

Corollary 2.9 gives \(a_-+a_+<0\) under the same assumption.
Write \(A_\pm=\arctan(2a_\pm/\sqrt{3\rho})\).
Since \(a_-<-a_+\), strict monotonicity and oddness
of \(\arctan\) give \(A_-<-A_+\). Thus the finite model strip
for the scaled metric has ordered endpoints
\(A_-/2,-A_+/2\) in \((-\pi/4,\pi/4)\).
Its end means agree with the scaled lower bounds in the
outward divergence convention of Remark 1.4. Theorem 2.7
bounds \(\sqrt{\rho/12}\,d\) above by the model width,
giving \eqref{eq:mean}
\cite[Remark 1.4, Theorems 2.7--2.8, Corollaries 2.9 and 2.25, and Lemma 6.3]{Raede2022Bands}.
\end{proof}

The non-strict finite-end inequality permits equality.
For example on a compact interval in
\((-\pi/(2\kappa),\pi/(2\kappa))\) take
\(dt^2+\cos(\kappa t)h_0\), where \(h_0\) is flat on \(T^3\)
and \(\kappa=\sqrt{\rho/3}\).
Its scalar curvature is \(3\kappa^2\).
The two outward means are
\((3\kappa/2)\tan(\kappa t_-)\) and
\(-(3\kappa/2)\tan(\kappa t_+)\), and its width is \(t_+-t_-\).
Substitution gives equality in \eqref{eq:mean}. The independent
strict obstruction \eqref{eq:width} includes equality on
its recognizing side. On a unit-volume positive Yamabe
representative, the constant scalar curvature is \(Y([g])\),
so the width criterion becomes \(Y([g])d^2\ge3\pi^2\).
The distance and scalar lower bound concern the metric on
the marked band itself.

\begin{proposition}[A slice or a stable separator]\label{prop:slice}
Suppose \(g=f^2dt^2+h_t\), \(f>0\), in the given marking,
and write \(K=(2f)^{-1}\partial_t h\), \(H=\tr_h K\), and
\(\Delta_h=\divg_h\nabla_h\). The slice is PSC-dominated
if
\[
 R_g+H^2+|K|^2+2f^{-1}\partial_t H
       +\tfrac32|\nabla_h\log f|^2
\]
is nonnegative everywhere and positive somewhere on the
slice. Alternatively, a smooth closed two-sided stable
minimal component \(S\) in the band with \(R_g>0\) along
it and nonzero projection degree certifies the slice.
\end{proposition}
\begin{proof}
The Gauss and normal mean curvature variation equations
with normal \(f^{-1}\partial_t\) give
\[
 R_h=R_g+H^2+|K|^2+2f^{-1}\partial_tH
                +2f^{-1}\Delta_h f.
\]
Then \(v=f^{1/4}\) is positive and
\[
 v^{-1}(-8\Delta_h+R_h)v
 =R_g+H^2+|K|^2+2f^{-1}\partial_tH
                  +\tfrac32|\nabla_h\log f|^2 .
\]
Pairing \((-8\Delta_h+R_h)v\) with the positive first
eigenfunction proves that the first eigenvalue is positive.
The conformal metric from that eigenfunction has positive
scalar curvature on the slice.

For the minimal component use stability for every smooth
normal variation, with no constraint on its integral.
The second variation with the Gauss equation gives
\[
 \int_S(2|\nabla u|^2+R_Su^2)\ge
 \int_S(R_g+|K|^2)u^2 .
\]
As \(S\) is compact and \(R_g>0\) there,
\(-8\Delta_S+R_S\) has positive first eigenvalue.
Projection is then the required nonzero-degree map from
a positive metric on \(S\). For a collection representing
the slice class, some component has nonzero degree.
\end{proof}

The width, slice and separator implications give direct
smooth recognition or reduction through the marked cuts.
They require no energy transfer between distinct metrics.
An application to a cap on the original metric needs an
embedded comparison carrying the exterior collars and the
actual gluing, as well as the separate chiral estimates
below \cite{Baker2026CollarDeformation,Baker2026CollarProgram}.

\section{Spherical cuts on a physical source}\label{sec:budget}
For a subset \(A\subset M\) and chirality \(\sigma\in\{+,-\}\)
put
\[
 E_\sigma(g;A)=\int_A|W_g^\sigma|_{\mathrm{HS}}^2\,dV_g.
\]
These are the Hilbert--Schmidt norms of the real Weyl
operators on \(\Lambda^\sigma\). Use the orientation of
\(M\) on all regions and caps. On a closed homotopy sphere
the two total energies agree by signature.
Set \(T=4\pi^2\) and
\[
 I=I^+(M)=\inf_{Y([g])>0} E_+(g;M),
\]
with the empty infimum \(+\infty\).
An actual positive class with \(E_+\le T\) recognizes
\(S^4\). An inclusive infimum criterion follows from the
universal non-effective gap of the master result
\cite[Recognition and uniform recognition by the Weyl spectrum]{Baker2026ChiralMaster}
\cite{ChenZhu2014Conformal,Fu2018FourManifolds}.

\subsection{From the joined peak to two positive caps}
Suppose a physical marking
\(\Phi:S^3\times(-w_L,w_R)\to M\) embeds a bicollar and
identifies its two complete sides \(P_L,P_R\).
A metric comparison must give the native widths, the
angular mark and its inverse, and all changes on each
side, including any connectors. Let a smooth joined
path starting at \(g\) with \(Y([g])>0\) have strict peak
\[
 H=\sup_s E_+(g_s;M)<2T
\]
and end at \(h\) with a round product germ in that
marking up to a smooth positive conformal factor.
The sign of the conformal Laplacian cannot change:
a zero first eigenvalue supplies a smooth scalar-flat
conformal representative \(q\) by its positive
eigenfunction. Signature and Gauss--Bonnet give
\[
 E_+(q;M)=2T+\tfrac14\int_M|\Ric_q|^2\,dV_q .
\]
Conformal invariance contradicts \(H<2T\).
For a unit-volume positive Yamabe representative
\(\widetilde g_s\), the same identities give
\[
 Y([g_s])^2
 =48(2T-E_+(g_s;M))
      +12\int_M|\Ric^\circ_{\widetilde g_s}|^2\,dV_{\widetilde g_s}
 \ \ge\ 48(2T-H).
\]
Here the Yamabe theorem supplies the smooth constant
scalar curvature representative
\cite{Schoen1984Yamabe,LeeParker1987Yamabe}
\cite[scalar-flat wall and short marked cut]{Baker2026ChiralMaster}.

Choose the conformal representative exact on a smaller
round product and stretch that product to arbitrary
finite length. The exact portion has zero Weyl energy.
Each extension of the path on the same marked joined
manifold therefore gives the same lower bound
\(y_0=\sqrt{48(2T-H)}>0\) for the Yamabe constant.
A compactly supported Yamabe test on a complete
retained side with the attached half-infinite cylinder
extends by zero to a long enough joined stretch. Hence
the cylindrical side has the same lower bound for
compactly supported tests. Conformally compactify the
round end at \(r=e^{-u}\) for cylinder parameter \(u\)
and multiply by \(e^{-2u}\) near infinity. This is
the smooth flat ball metric there in the actual
angular marking. Conformal covariance and the zero
\(H^1\) capacity of the added point in dimension four
give \(Y\ge y_0\) on the resulting closed capped side.
Thus for either \(j\in\{L,R\}\) and chirality the
capped energy is exactly \(E_\sigma(h;P_j)\).
In particular there are positive caps on \emph{both}
sides of the same embedded sphere
\cite[short marked cut and exterior replacement]{Baker2026ChiralMaster}.

\subsection{Physical reassembly and finite inequalities}
An embedded \(S^3\) in a homotopy four-sphere separates,
since a nonseparating hypersurface would give a loop
with algebraic intersection one. Van Kampen and
Mayer--Vietoris imply that capping both complete sides
by marked standard balls gives simply connected
homology spheres. They are homotopy spheres by the
Hurewicz and homological Whitehead theorems.
The marked diffeomorphisms along their cap boundaries
are the ones from \(\Phi\) with the opposite outward
normals. If one capped side is standard, Palais's
isotopy extension for a parametrized cap ball makes
its complement a standard ball. Hatcher's \(S^3\)
diffeomorphism theorem extends the resulting boundary
map across that ball. Original reassembly is therefore
oriented-diffeomorphic to the other capped side,
relative to its preserved exterior outside the
gluing collar
\cite[Theorem B]{Palais1960Extending}
\cite[p.~553]{Hatcher1983Smale}
\cite[finite marked cut criterion]{Baker2026ChiralMaster}.

For clarity, write \(m_j^\sigma=E_\sigma(g;P_j)\),
\(e=E_+(g;M)\), and let \(\gamma_j^\sigma\ge0\) be
an upper bound on the surcharge of a positive cap
over \(m_j^\sigma\). On either sign \(m_L^\sigma+
m_R^\sigma=e\). If the cap energy sum is at most \(2T\),
in particular if \(e+\gamma_L^\sigma+\gamma_R^\sigma\le2T\),
at most one cap can have energy greater than \(T\).
If both are at most \(T\), the original source is
standard. If one exceeds \(T\), the opposite cap is standard,
and the original source has the smooth type of the unique cap
of energy greater than \(T\). On a nonstandard fixed type this forces,
with \(d=e-I\),
\begin{equation}\label{eq:nonsplit}
 \min(m_L^\sigma,m_R^\sigma)
 \le d+\max(\gamma_L^\sigma,\gamma_R^\sigma).
\end{equation}
Indeed the actual surviving cap is a positive
competitor for \(I\). The cap sum test is inclusive
after positivity has been established. The joined
path argument for that positivity uses a strict peak.
Cap positivity certified independently also permits
\eqref{eq:nonsplit}.

\subsection{The original charged disk and its exterior}
Take \(I<2T\) on a hypothetical nonstandard type and
a minimizing sequence of smooth unit-volume positive
Yamabe metrics \(g_i\), with equal total chiral energy
\(e_i\to I\) and \(d_i=e_i-I\). The master charged-disk
theorem supplies a subsequence, parametrized standard
disks \(D_i\) with marked bicollars, a fixed chirality
\(\sigma_*\), and constants \(a,b>0\) such that
\begin{equation}\label{eq:charge}
 m_i^{\sigma_*}=E_{\sigma_*}(g_i;D_i)\ge a,\qquad
 v_i^{\sigma_*}=E_{\sigma_*}(g_i;V_i)\ge b,\qquad
 V_i=M\setminus\intt D_i .
\end{equation}
Both masses are on the complete sides, including the
original parts of the collar. In each sign set
\(m_i^\sigma=E_\sigma(g_i;D_i)\) and
\(v_i^\sigma=E_\sigma(g_i;V_i)=e_i-m_i^\sigma\)
\cite[charged disk throughout the subcritical band]{Baker2026ChiralMaster}.

Fix one index with strict headroom. If a path to a
matched product germ has joined peak \(H_i<2T\),
its capped exterior is on the original smooth type:
the actual disk parametrization extends the angular
mark over the removed disk, and the round radial
compactification extends the same angular mark over
the new ball. The resulting diffeomorphism is the
identity on the retained exterior away from its
gluing collar. Writing \(h_i\) for the endpoint
and \(\Delta_{V,i}^\sigma=E_\sigma(h_i;V_i)-v_i^\sigma\),
the capped exterior has positive Yamabe constant
and satisfies
\begin{align}
 E_\sigma(\widehat g_i;M)
 &=e_i-m_i^\sigma+\Delta_{V,i}^\sigma\ \ge I,\label{eq:ext}\\
 m_i^\sigma&\le d_i+\Delta_{V,i}^\sigma,\label{eq:mass}\\
 \Delta_{V,i}^+-\Delta_{V,i}^-&=m_i^+-m_i^-.\label{eq:asym}
\end{align}
The last identity follows by signature on the capped
exterior. The full signed exterior change includes
every portion of the exterior modified by the path.
A certified upper bound \(G_i^\sigma\) for that change
gives a contradiction if \(m_i^\sigma-G_i^\sigma>d_i\).
If \(e_i-m_i^\sigma+G_i^\sigma\le T\), the positive
competitor gives direct recognition. The complete
disk-side change enters the joined peak even though
the competitor retains the exterior
\cite[exterior replacement]{Baker2026ChiralMaster}.

For example, suppose the physical annular splice
writes each changed chiral density as a factor of
absolute value at most one times the raw amplitude
plus a remainder with \(L^2\) norm at most \(\eta\)
in the fixed reference measure. If the full changed
strip has raw mass \(c\), expansion of the squared
\(L^2\) norm gives an upper increment
\(2\sqrt c\,\eta+\eta^2\). The estimate can be applied
with separate remainder bounds on the whole strip and
on its exterior half, using the corresponding raw
masses. A joined increment bound \(P_i^\sigma\) valid
along the whole path gives the strict wall test
\(e_i+\min_\sigma P_i^\sigma<2T\), because the two
joined totals agree. Separate exterior bounds still
have to be compared to their own signed disk masses.
These estimates require the supplying metric in the
physical chart and the actual widths used in the cutoff
\cite[first-order splice]{Baker2026ChiralMaster}
\cite{Baker2026FixedDescent,Baker2026CollarDeformation}.

\begin{theorem}[Exact face obstruction at one index]\label{thm:exact}
Fix a physical index with \(e_i<2T\).
Suppose the boundary of \(D_i\) is umbilic and its induced
metric is conformal to the unit round metric in the actual
disk marking. Then
\[
 E_+(g_i;D_i)=E_-(g_i;D_i)\le d_i .
\]
Consequently this exact face is excluded on every
sufficiently late selected disk in \eqref{eq:charge}.
\end{theorem}
\begin{proof}
Choose the normal toward \(V_i\), with second fundamental
form \(K\) and trace \(H_0\) at the face. Under change
to \(e^{2f}g_i\) the second fundamental form is
\(e^f(K+\nu(f)h_0)\). Prescribe \(f\) to make the face
metric \(\rd\), and choose \(\nu(f)=-H_0/3\) since
\(K^\circ=0\). Extend these jets smoothly with a
cutoff in the marked bicollar. The new normal flow
extends the same angular marking and gives a Gaussian
chart \(dt^2+h_t\) with \(h_0=\rd\) and
\(\partial_t h_0=0\).

Within this fixed chart choose arbitrarily small
positive native subwidths \(w_D,w_V\). Let
\(\psi(t)=\chi(-t/w_D)\) on the disk side and
\(\chi(t/w_V)\) on the exterior side, with a smooth
\(\chi:[0,1]\to[0,1]\) constant one near zero and zero near one.
The smooth path replaces \(h_t\) by
\[
 \rd+(1-s\psi(t))(h_t-\rd),\qquad 0\le s\le1 .
\]
It is positive by convexity and agrees with the
source near both outer faces. Its endpoint is a
round product in the central strip. Multiplying
back by \(e^{-2f}\) restores the original starting
metric and preserves all the chiral energy integrals.
In a fixed finite cover of the chart, derivatives of
the mismatch through order two in the normal variable
are bounded by \(O_i(w^2),O_i(w),O_i(1)\) on a strip
of width \(w\), with corresponding tangential bounds
where needed. Thus all differentiated cutoff terms
needed for curvature stay bounded uniformly in
\(s,w_D,w_V\) at this source. The inverse metrics and
volume densities stay bounded as well. Integrating
the uniform curvature bound over the supports yields
\[
 H_i\le e_i+C_i(w_D+w_V),\qquad
 |\Delta_{V,i}^\sigma|\le C_i w_V,\qquad
 |E_\sigma(h_i;D_i)-m_i^\sigma|\le C_i w_D .
\]
The fixed-source constants and the refined leading
coefficients belong to the companion deformation
calculation \cite{Baker2026CollarDeformation}.

Choose positive widths within the strict energy
headroom at this index. They give actual positive
capped metrics and \eqref{eq:mass}--\eqref{eq:asym}.
Let the two widths tend to zero at this fixed
source. Equation~\eqref{eq:asym} gives equality of
the raw disk masses, and \eqref{eq:mass} gives the
inclusive bound by \(d_i\) in both signs. In
particular if a mass exceeds \(d_i\), a sufficiently
small positive exterior width already gives a
strict finite contradiction. Finally \(d_i\to0\)
and \(m_i^{\sigma_*}\ge a\) on the selected
subsequence, so every sufficiently late index
violates the exact-face bound. The width limit
at one index precedes the passage along the
sequence, and its constants may depend on \(i\).
\end{proof}

At an attained subcritical positive minimizer the
same exact face condition forces zero Weyl mass on
the whole disk in both chiralities. For a sequence
of approximate faces the compactness argument in the
companion requires physical collars smoothly
precompact up to smooth conformal gauge on a common
slab in the marking, with positive definite limits
and enough room for each permitted splice. On such
a precompact subsequence of the charged disks it
implies a positive lower limit for the integrated
Cotton and umbilicity defect
\cite{Baker2026CollarDeformation}.

\subsection{Balanced enlargements with a protected core}
Fix an integer \(k\ge1\). The master balanced
selection preserves an outer disk \(Q_i\) containing
the original charged disk and its marked bicollar.
It produces a larger parametrized standard disk
\(B_i\) containing \(Q_i\).
Writing \(V_i^B=M\setminus\intt B_i\),
\(q_i^\sigma=E_\sigma(g_i;Q_i)\), and
\(b_i^\sigma=E_\sigma(g_i;B_i)\), its division is
\begin{equation}\label{eq:balanced}
 \begin{aligned}
 E_\sigma(g_i;V_i^B)&=2^{-k}(e_i-q_i^\sigma)+r_i^\sigma,\\
 b_i^\sigma&=(1-2^{-k})e_i+2^{-k}q_i^\sigma-r_i^\sigma,
 \qquad r_i^\sigma\longrightarrow0 .
 \end{aligned}
\end{equation}
The construction uses an auxiliary ball containing
all but a small tolerance of each of the four
exterior measures, the two chiral measures, the
spectral measure, and the original full-curvature
measure. One affine hyperplane in its chart
bisects all four. A rounded half joined by a thin
tube gives a boundary connected sum of standard
disks while preserving the protected \(Q_i\).
Iteration at fixed depth gives \eqref{eq:balanced}.
This provides thin strips of small raw mass, with
a separate approximate matching question for the
selected boundary
\cite[balanced disk-side saving]{Baker2026ChiralMaster}
\cite{StoneTukey1942Sandwich}.

For any matched path with joined peak \(H_i<2T\),
put \(e_i'=E_\sigma(h_i;M)\),
\(S_i^\sigma=b_i^\sigma-E_\sigma(h_i;B_i)\),
and
\(\Delta_{V^B,i}^\sigma
   =E_\sigma(h_i;V_i^B)-E_\sigma(g_i;V_i^B)\).
The positive exterior cap is again on \(M\),
so \eqref{eq:ext} and \(e_i'\le H_i\) give
\begin{align}
 \Delta_{V^B,i}^\sigma&\ge b_i^\sigma-d_i,\label{eq:bigext}\\
 S_i^\sigma&\ge b_i^\sigma+I-H_i
                   > b_i^\sigma+I-2T,\label{eq:save}\\
 \Delta_{V^B,i}^\sigma&=S_i^\sigma+e_i'-e_i
                   \ge S_i^\sigma-d_i.\label{eq:couple}
\end{align}
The last inequality uses that \(h_i\) on the joined
source is itself positive. Substituting
\eqref{eq:balanced} into the strict wall bound in
\eqref{eq:save} retains the term \(2^{-k}q_i^\sigma\).
At \(k=1\) and \(I=4T/3\) the limiting lower bound
on the required savings in \(\sigma_*\) is at least
\(a/2\). The retained raw exterior at \(k=1\)
is eventually below \(T\) in both signs since
\(I<2T\). Its positive exterior cap nevertheless
has energy at least \(I\) on the original type.

The master uniform recognition gap gives
\(I>T+\varepsilon_*\) on a nonstandard type for a
universal \(\varepsilon_*>0\). Choose one fixed
finite \(k\) with \(2^{-k}(T+\varepsilon_*)\le
\varepsilon_*\). Then the lower limit of the
strict wall bound in both signs is at least
\((2-2^{-k})I-2T>\varepsilon_*\).
A change confined to a final seam strip on the
disk side saves at most the raw mass on that strip.
Consequently strips whose raw masses tend to zero
cannot produce a matched path with the strict
peak at sufficiently late indices for this \(k\).
If the full core \(Q_i\) is preserved up to conformal
gauge by a larger deformation, the disk endpoint
mass in each sign is at least \(q_i^\sigma\), so
\[
 H_i\ge I+\max_\sigma q_i^\sigma .
\]
At a tested index a core mass of at least \(2T-I\)
in either sign rules out this strict peak with that
core preserved. Equations~\eqref{eq:bigext}--\eqref{eq:couple}
specify the exterior cost and the complete disk-side
saving for a coupled deformation on a given source.

\section{The next unrecognized presentations}\label{sec:open}
A winding-one knot with three geometric crossings is
the first mixed candidate beyond a single geometric
intersection. The bigon and peripheral arguments require
a chosen disk or control of a particular meridian.
A missing bigon, an obstructed peripheral element or a
failed scalar test in one order leaves the other marked
cuts and slides available. The unresolved mixed class
after these tests lacks a verified reducing move,
belt-circle recognition or PSC-dominated intermediate
boundary in the tested forward and dual diagrams.
A metric certificate here needs a physical comparison
to one of those levels, for example an actual compact
band or a stable separator of nonzero projection degree.

For one lower 1-handle and more middle handles, the
winding-one move preserves all remaining movable tubes
and strictly reduces \(n_2\). The first unresolved
three-middle-handle successor after that move is the
actual two-component link when its second prefix and
other one-sided recognition tests are undecided.
A general one-sided three-component link and a presentation
with several lower 1-handles require additional link
or trace selection. A result on a freely selected
standard-type control supplies the explicit metric and
regluing on that control only.

The charged-disk comparison is on the original smooth
minimizing source. Its remaining cases include positive
face defect, loss of compactness in the physical markings
or conformal gauges, and absence of a measured coupled
deformation whose joined peak and full exterior cost pass
the finite comparison. An identification of a handle cut
with such a disk requires a smooth marked comparison
and the same physical metric. Finally, arbitrary homotopy
four-spheres supply Morse presentations, while the
variational comparison starts only after a positive class
of energy below \(8\pi^2\) has been supplied. The low-energy
immersed-height route has a distinct input, an immersion
with induced energy below the bound of its own theorem.
A direct topological reduction requires neither entry
hypothesis \cite{Baker2026MarkedHeights,Baker2026ChiralMaster,Baker2026CollarProgram}.

\begingroup
\raggedright
\bibliographystyle{unsrturl}
\bibliography{references}
\endgroup
\end{document}
